%----------------------------------------------------------------------------------------------
%-----------------------------------------     5. An1a     ------------------------------------
%----------------------------------------------------------------------------------------------
\section{Wichtiges aus An1a}


%----------------------------------------------------------------------------------------------
%---------------------------     5.1 Gebrochenratinale Funktionen     -------------------------
%----------------------------------------------------------------------------------------------
\subsection{Gebrochenratinale Funktionen}{63m 67}

\[
f(x) = \frac{P_n(x)}{Q_m(x)} = \frac{a_nx^n + a_{n-1}x^{n-1} + \ldots + a_1x^1 + a_0}{b_mx^m + b_{m-1}x^{m-1} + \ldots + b_1x^1 + b_0}
\]

\begin{tabular}{l l}
    Echt gebrochen      &   $n < m$ \\
    gleich Gradig       &   $n = m$ \\
    Unecht gebrochen    &   $n > m$
\end{tabular}

\vspace{0.2cm}
Jede unecht gebrochene Funktion lässt sich als Summe einer ganzrationalen Funktion und einer echt gebrochenen Funktion schreiben. $\Rightarrow$ \textbf{Polynomdivision}


%----------------------------------------------------------------------------------------------
%----------------------------------     5.2 Grenzwerte     ------------------------------------
%----------------------------------------------------------------------------------------------
\subsection{Grenzwerte}{471}

Die Folgen ${\textcolor{red}{a_n}}$ und ${\textcolor{blue}{b_n}}$ sind konvergent!\vspace{0.5em}
\\
\renewcommand{\arraystretch}{2}
\begin{tabular}{ll}
    $\lim\limits_{n \to \infty} \abs{\textcolor{red}{a_n}}$ \quad \;\;\; $=$   $\abs{\lim\limits_{n \to \infty} \textcolor{red}{a_n}} $\\

    $\lim\limits_{n \to \infty} \left(\textcolor{red}{a_n} + \textcolor{blue}{b_n}\right)$  $=$   $\lim\limits_{n \to \infty} \textcolor{red}{a_n} + \lim\limits_{n \to \infty} \textcolor{blue}{b_n}$ &
    $\lim\limits_{n \to \infty} \left(\textcolor{red}{a_n} - \textcolor{blue}{b_n}\right)$  $=$   $\lim\limits_{n \to \infty} \textcolor{red}{a_n} - \lim\limits_{n \to \infty} \textcolor{blue}{b_n}$\\
    
    $\lim\limits_{n \to \infty} \left(\textcolor{red}{a_n} \cdot \textcolor{blue}{b_n}\right)$ \; $=$   $\lim\limits_{n \to \infty} \textcolor{red}{a_n} \cdot \lim\limits_{n \to \infty} \textcolor{blue}{b_n}$ &
    $\lim\limits_{n \to \infty} \left(\frac{\textcolor{red}{a_n}}{\textcolor{blue}{b_n}}\right)$ \quad \;\:\, $=$   $\frac{\lim\limits_{n \to \infty} \textcolor{red}{a_n}}{\lim\limits_{n \to \infty} \textcolor{blue}{b_n}}$, \quad für $\lim\limits_{n \to \infty} \textcolor{blue}{b_n} \neq 0$ 
\end{tabular}


%----------------------------------------------------------------------------------------------
%--------------------------     5.3 Grenzwerte gegen unendlich     ----------------------------
%----------------------------------------------------------------------------------------------
\subsection{Grenzwerte \texorpdfstring{$\longrightarrow \infty$}{→ ∞}}
\[
    \lim\limits_{n \to \infty}\left(\frac{-2n^2+4n-5}{8n^2-3n+7}\right)
\]
\begin{tabular}{lcl}
    $ \displaystyle{f(n)= \frac{-2n^2+4n-5}{8n^2-3n+7}}$                    & Algebraisch erweitern mit $\frac{1}{n^2}$: & $\displaystyle{f(n)=\frac{-2+\frac{4}{n}-\frac{5}{n^2}}{8-\frac{3}{n}+\frac{7}{n^2}}}$\\
    $ \displaystyle{f(n)=\frac{-2+\frac{4}{n}-\frac{5}{n^2}}{8-\frac{3}{n}+\frac{7}{n^2}}}$ & $\xrightarrow{n \to \infty}$               & $\displaystyle{f(n)=\frac{-2}{8} = -\frac{1}{4}}$
\end{tabular}\\

\[
    \lim\limits_{n \to \infty}\left(\frac{-2n^2(4n-5)}{8n^2(-3n+7)}\right)
\]
\begin{tabular}{lcl}
    $ \displaystyle{f(n)= \frac{-2n^2(4n-5)}{8n^2(-3n+7)}}$                 & Algebraisch erweitern mit $\frac{1}{n^3}$: & $\displaystyle{f(n)=\frac{-8+\frac{10}{n}}{-24+\frac{56}{n}}}$\\
    $ \displaystyle{f(n)=\frac{-8+\frac{10}{n}}{-24+\frac{56}{n}}}$         & $\xrightarrow{n \to \infty}$               & $\displaystyle{f(n)=\frac{-8}{-24} = \frac{1}{3}}$ 
\end{tabular}


%----------------------------------------------------------------------------------------------
%----------------------------------     5.4 Grenzwerte     ------------------------------------
%----------------------------------------------------------------------------------------------
\subsection{Rechnen \texorpdfstring{$\Longrightarrow \infty$}{→ ∞}}

\subsubsection{Bestimmte Formen}

\renewcommand{\arraystretch}{2}
\begin{tabular}{lll}
    $\infty + \infty = \infty$      & $-\infty - \infty = -\infty$          & $0 \cdot [a,b] = 0 \cdot \text{beschr"ankt} = 0$  \\
    $g + \infty = \infty$           & $g - \infty = -\infty$                & 	$(g \in \mathbb{R})$                            \\
    $\infty \cdot \infty = \infty$  &	$-\infty \cdot (\infty) = -\infty$  &
    $g \cdot \infty =
    \begin{cases}
        \infty   & g > 0 \\ 
        -\infty & g < 0 
    \end{cases}$        \\ 		
    $\frac{1}{\infty} = 0$          &	$\frac{g}{\infty} = 0$              & 	$g \in \mathbb{R}$	                            \\
    $\frac{\infty}{0+} = \infty$	& $\frac{\infty}{0-} = -\infty$
    & $\frac{\infty}{g} =
    \begin{cases}			
        \infty & g > 0 \\
        -\infty & g < 0
    \end{cases} $ \\
    $\frac{1}{0+} = \infty$         & $\frac{1}{0-} = -\infty$	            & $g \in \mathbb{R} - {0}$                          \\
\end{tabular}

\renewcommand{\arraystretch}{1}


%----------------------------------------------------------------------------------------------
%------------------------------------     5.5 Stetigkeit     ----------------------------------
%----------------------------------------------------------------------------------------------
\subsection{Stetigkeit}{59-63, 127}

$$ \text{Definition Stetigkeit:} \quad \lim \limits_{x \to x_0} f(x) = f(x_0) $$
Eine Funktion ist stetig, wenn der Funktionsgraph gezeichnet werden kann, ohne dass der Stift abgesetzt werden muss.

\begin{center}
	\includegraphics[width=0.98\columnwidth]{images/stetigkeit.PNG} 
\end{center}


%----------------------------------------------------------------------------------------------
%--------------------------------     5.6 Wichtige funktion     -------------------------------
%----------------------------------------------------------------------------------------------
\subsection{Wichtige Funktionen}

\subsubsection{Exponentialfunktion}{73}
\[
\e^1 = \lim\limits_{n \to \infty} \left(1+\frac{1}{n} \right)^n \quad \quad \e^{-1} = \lim\limits_{n \to \infty} \left(1-\frac{1}{n} \right)^n \quad \quad \e^x = \lim\limits_{n \to \infty} \left(1+\frac{x}{n} \right)^n
\]
\begin{minipage}{0.45\columnwidth}
    \includegraphics[width=0.95\columnwidth]{images/exp-funktion.PNG}
\end{minipage}
\hfill
\begin{minipage}{0.5\columnwidth}
    \raggedright
    \textbf{Definitions- / Wertebereich:}
    $D_f = \mathbb{R} \rightarrow W_f = \mathbb{R^+}$ \\
    \vspace{0.4cm}
    
    \textbf{Einschliessung:}
    
    \renewcommand{\arraystretch}{1.6}
    \begin{tabular}{ll}
        $\e^x \geq 1 + x$ & f"ur $x \in \mathbb{R}$ \\
        $\e^x \leq \frac{1}{1-x}$ & f"ur $x < 1$  
    \end{tabular}
\end{minipage}
\hrule

\subsubsection{Hyperbolische Funktionen}{89}

$$ \e^x = \frac{1}{2} (\e^x - \e^{-x}) + \frac{1}{2} (\e^x + \e^{-x}) = \sinh(x) + \cosh(x) $$ 

\renewcommand{\arraystretch}{1.5}
\begin{tabular}{ll}
    $\sinh(x) = \frac{1}{2} (\e^x - \e^{-x})$ & $\mathbb{R} \to \mathbb{R}$  \\
    $\cosh(x) = \frac{1}{2} (\e^x + \e^{-x})$ & $\mathbb{R} \to [1; \infty)$ \\
    $\tanh(x)$ = $\frac{\sinh(x)}{\cosh(x)} = \frac{\frac{1}{2} (\e^x - \e^{-x})}{\frac{1}{2} (\e^x + \e^{-x})}$ & $\mathbb{R} \to (-1; 1)$ \\	

    $\vert \sinh(x) \vert < \cosh(x)$	& 
\end{tabular}
\renewcommand{\arraystretch}{1}
